% Insert as A.68.  Fixed abelian-cover quotients retain full Tate modules.
\begin{statement}{A.68}{Proved}{Ordinary ghosts in every fixed full-Tate abelian-cover quotient}
Let $\Pi=\pi_1(U_{\overline\Q})$, put
$A=\Pi^{\rm ab}=\widehat\Z(1)^2$, and fix the point-$3$ section $s_0$.
For $n\geq1$ set
\[
 H_n=\ker(\Pi\longrightarrow A/nA),\qquad
 H_n'=\overline{[H_n,H_n]},\qquad E^{[n]}=\pi_1(U)/H_n'.
 \tag{A.68.1}
\]
These are actual arithmetic quotients of the full geometrically
metabelian quotient, because $\Pi''\subset H_n'$.
Every abelian section whose class relative to $s_0$ belongs to
$n^2H^1_{\rm cts}(G_\Q,A)$ lifts to an actual section of $E^{[n]}$.

Consequently, for each fixed $n$, there are ordinary infinite-support
unit-equation pairs lifting to $E^{[n]}$.
This quotient retains the entire adelic Tate module of the generalized
Jacobian of the degree-$n^2$ cover below, including its compact
Jacobian part.  All coefficient primes and all torsion depths in this
one quotient are retained.

The full geometrically metabelian extension is the inverse limit of
these quotients.  The proof specifies the transition maps and the
continuous obstruction to passing from separate lifts to a compatible
lift; no full metabelian lift of an ordinary ghost is asserted.
\end{statement}
\begin{proof}
\emph{The pointed covering and its actual abelianization.}
Let
\[
 V_n:\quad 3u^n-2v^n=1,\quad uv\ne0,\qquad
 P_n=(1,1),\qquad f_n(u,v)=3u^n .
 \tag{A.68.2}
\]
The map $f_n:V_n\to U$ is finite \'etale of degree $n^2$ and takes
$P_n$ to $3$.  Over $\overline\Q$ its deck group is $\mu_n^2$, and
its fundamental group is precisely $H_n$.
The rational point $P_n$ specifies the same arithmetic splitting
$s_0$ on this open subgroup.  Its smooth compactification is
\[
             C_n:\quad3X^n-2Y^n=Z^n.
 \tag{A.68.3}
\]
It is geometrically smooth; the three boundary divisors
$D_{0,n}=(X=0)$, $D_{1,n}=(Y=0)$ and $D_{\infty,n}=(Z=0)$
are disjoint reduced finite \'etale schemes of degree $n$ over $\Q$.
Write $D_n$ for their union and $J_n$ for the generalized Jacobian
of $(C_n,D_n)$, based by the Abel map at $P_n$.
Its exact sequence is
\[
 1\longrightarrow
 T_n=\bigl(\operatorname{Res}_{D_n/\Q}\G_m\bigr)/\G_m
 \longrightarrow J_n\longrightarrow\operatorname{Jac}(C_n)
 \longrightarrow0 .
 \tag{A.68.4}
\]
Over $\overline\Q$ the character lattice of $T_n$ is the degree-zero
divisors supported on $D_n$, and its cocharacter lattice is
$\Z[D_n(\overline\Q)]/\Z\sum_{P\in D_n(\overline\Q)}e_P$.
Both descriptions retain the natural Galois permutation action.

The Abel map gives the actual equivariant identification
\[
                H_n^{\rm ab}\ \simeq\ T J_n .
 \tag{A.68.5}
\]
For the fundamental-group comparison, one may use
\cite[Proposition A.8(ii),(iii), pp.~232--234]{BIR_MochizukiAlbanese2012}.
For part (ii), first extend scalars to a finite Galois field over which
every point of $D_n$ is rational.  Its geometrically irreducible boundary
hypothesis then holds, and the character group in that proposition is
the degree-zero subgroup of $\Z[D_n(\overline\Q)]$.
The identification is functorial in the pointed curve, so descent along
that Galois extension gives precisely the permutation action and the
$\Q$-torus in (A.68.4).  Part (iii) applies to the pointed smooth
curve over $\Q$ itself.
The torsion-free abelianization agrees with the full
abelianization here: $H_n$ is free profinite of rank $n^2+1$.
The semi-abelian fundamental group is its full adelic Tate module.
The character-lattice description in (A.68.4) also proves the torus
part directly from the divisors of units.

The pointed units $u,v$ induce a homomorphism
$\phi_n:J_n\to\G_m^2$.
Since $f_n/3=u^n$ and $(1-f_n)/(-2)=v^n$, the actual map induced
by the covering on geometric abelianizations is
\[
 \rho_n:T J_n\longrightarrow A,\qquad
                        \rho_n=n\,T\phi_n .
 \tag{A.68.6}
\]
Its image is $nA$, by the definition of $H_n$.

\emph{A rational boundary-torus map and the factor $n^2$.}
Define $j_n:\G_m^2\to T_n\subset J_n$ by assigning the value $a$
to the whole finite \'etale scheme $D_{0,n}$, the value $b$ to
$D_{1,n}$, and the value $1$ to $D_{\infty,n}$, and taking the
class modulo the diagonal $\G_m$.
This is a morphism over $\Q$, without choosing any individual cusp.
On cocharacters its columns are
\[
          \sum_{P\in D_{0,n}}e_P,\qquad
          \sum_{P\in D_{1,n}}e_P.
 \tag{A.68.7}
\]
Each sum is invariant under the full Galois permutation action.

The two character divisors are exactly
\[
             \operatorname{div}(u)=D_{0,n}-D_{\infty,n},
 \qquad       \operatorname{div}(v)=D_{1,n}-D_{\infty,n}.
 \tag{A.68.8}
\]
Pairing (A.68.8) with (A.68.7) gives the matrix $nI_2$.
Equivalently, the restriction of $\phi_n$ to the torus takes the
two indicated constant functions to their norms $a^n,b^n$.
The sign convention is fixed by positive boundary meridians:
the character of a unit on the torus is its divisor, as follows
also from Weil reciprocity for a principal divisor disjoint from
the boundary.
Thus the actual algebraic-group and Tate-module identities are
\[
       \phi_n j_n=[n],\qquad
       \rho_n\,t_n=n^2\operatorname{id}_A,\qquad
                          t_n=Tj_n:A\to TJ_n .
 \tag{A.68.9}
\]
This proves the multiplication factor and equivariance integrally.
The exact integer certificate
\path{code/metabelian_boundary_certificate.py} verifies the two
divisor/cocharacter matrices, their boundary permutations, and
the ramified tower identity (A.68.15); its saved output records
64 pairings, 162201 permutation checks and 1024 tower checks.

Represent the prescribed relative abelian class by $n^2a(g)$, where
$a\in Z^1_{\rm cts}(G_\Q,A)$.
The map
\[
                      s_n(g)=t_n(a(g))\,s_0(g)
 \tag{A.68.10}
\]
is an actual continuous section of
$(H_n/H_n')\rtimes_{s_0}G_\Q\subset E^{[n]}$.
Indeed $H_n/H_n'=TJ_n$ is abelian, $t_n$ is equivariant, and
the cocycle identity for $a$ is exactly the section identity.
Equation (A.68.9) makes its abelianization $n^2a$.
For an abelian section in the same cohomology class, lift an abelian
geometric conjugator and conjugate (A.68.10); hence the assertion is
also valid for the section-class image.

Choose $C$ in A.66 divisible by $24n^2$.
Its explicit convergent divided products make $A/3$ and $B/(-2)$
actual $n^2$th powers in the completed multiplicative group.
Equation (A.68.10) supplies the asserted section of $E^{[n]}$.
The ordinary local points and infinite support are those proved
in A.66.  The fixed valuations at two and three exclude cusp packets,
and infinite support excludes rational points.
No local point matching for the higher $TJ_n$ coordinates of
(A.68.10) is asserted.

\emph{The exact inverse system reaching the full metabelian quotient.}
For $n\mid m$, we have $H_m\subset H_n$ and $H_m'\subset H_n'$.
The transition homomorphism $E^{[m]}\to E^{[n]}$ is therefore the
actual quotient map.  Moreover,
\[
          \bigcap_{n\geq1}H_n'=\Pi'',\qquad
          \pi_1(U)/\Pi''\simeq\varprojlim_n E^{[n]} .
 \tag{A.68.11}
\]
Here is a direct proof of the first equality, including its topology.
The inclusion $\Pi''\subset H_n'$ holds for every $n$.
If $d\in\Pi'\setminus\Pi''$, a continuous finite abelian character
$\lambda$ of $\Pi'$ detects $d$.
Choose an open normal subgroup $W\subset\Pi$ with
$W\cap\Pi'\subset\ker\lambda$.
The rule $\widetilde\lambda(aw)=\lambda(a)$ on $\Pi'W$
is well-defined and is a homomorphism:
$[W,\Pi']\subset W\cap\Pi'$ ensures the required conjugation
invariance.  It is continuous.
The open subgroup $\Pi'W$ contains $H_n$ for some $n$, because
its image in $A$ is open.
Thus $H_n'\subset\ker\widetilde\lambda$, so $d\notin H_n'$.
An element outside $\Pi'$ is already outside $H_1'=\Pi'$.
This proves the equality.  The inverse-limit isomorphism follows
from compactness applied to the compatible cosets in $\pi_1(U)$.

Put
\[
 D=\Pi'/\Pi'',\qquad K_n=\Pi'/H_n'.
\]
The pointed cover gives the exact sequence and limit
\[
 0\longrightarrow K_n\longrightarrow TJ_n
              \xrightarrow{\rho_n}nA\longrightarrow0,
       \qquad D\simeq\varprojlim_nK_n .
 \tag{A.68.12}
\]
The underlying $\widehat\Z$ module $K_n$ is free of rank $n^2-1$.
Indeed $TJ_n=H_n^{\rm ab}$ has rank $n^2+1$, the displayed
map onto $nA$ is surjective, and $nA$ is free of rank two.
The transition maps $K_m\to K_n$ are surjective quotients of $D$.

For a fixed actual abelian section $s_{\rm ab}$, conjugation through
any of its lifts defines a well-defined action on the abelian group
$D$; denote the resulting module by $D_{s_{\rm ab}}$.
The quotient modules are $(K_n)_{s_{\rm ab}}$.
The ordered geometric lift $y^b x^a$ provides a continuous set
section over $s_{\rm ab}$, hence its actual obstruction is
\[
          o(s_{\rm ab})\in
             H^2_{\rm cts}(G_\Q,D_{s_{\rm ab}}).
 \tag{A.68.13}
\]
Its image in $H^2_{\rm cts}(G_\Q,(K_n)_{s_{\rm ab}})$ is the
obstruction to a lift to $E^{[n]}$.
On the fixed neighborhood where the relative abelian cocycle takes
values in $nA$, the action on $K_n$ is the point-$P_n$ action,
and this latter obstruction is exactly the connecting class for
(A.68.12).  Indeed conjugation by $H_n$ is trivial on $H_n/H_n'$,
so the pulled-back extension is the indicated extension of
abelian modules, with its actual cocycle map to $nA$.
This is the precise generalized-Jacobian comparison on that
neighborhood.

Taking the cofinal sequence $n=j!$, the continuous cochain complexes
give the exact sequence
\[
 \begin{split}
 0\longrightarrow&
 \varprojlim\nolimits^{1}_{j}
 H^1_{\rm cts}(G_\Q,(K_{j!})_{s_{\rm ab}})
 \longrightarrow H^2_{\rm cts}(G_\Q,D_{s_{\rm ab}})\\
 &\longrightarrow
 \varprojlim_jH^2_{\rm cts}(G_\Q,(K_{j!})_{s_{\rm ab}})
 \longrightarrow0 .
 \end{split}
 \tag{A.68.14}
\]
For this comparison, the cochain complexes really are degreewise
surjective inverse systems: each transition of the finite-rank free
compact $\widehat\Z$ modules has a continuous
$\widehat\Z$-linear set section.
It need not be equivariant to lift continuous cochains.
Thus their derived inverse limit is the ordinary inverse-limit
cochain complex, and its Milnor sequence is (A.68.14).
In particular separate vanishing of every quotient obstruction
would still leave precisely the displayed compatibility term.
Equivalently, a compatible sequence of actual framed sections of
the $E^{[j!]}$ is exactly an actual section of the inverse limit;
separate conjugacy classes are not substituted for such a sequence.

For clarity the toric maps used above have explicit transitions too.
If $m=kn$, the map $V_m\to V_n$ is $(u,v)\mapsto(u^k,v^k)$.
Every boundary cusp has ramification index $k$, and every boundary
cusp downstairs has $k$ preimages in its respective boundary set.
Hence the induced generalized-Jacobian Tate map $\tau_{m,n}$ satisfies
\[
                         \tau_{m,n}t_m=k^2t_n.
 \tag{A.68.15}
\]
This follows either by pushing the boundary meridians, each by its
ramification index, or directly from (A.68.7).
Accordingly the toric sections (A.68.10) are compatible wherever
one fixed abelian cocycle is divisible by all the required squares
and the divided cocycles are chosen compatibly.
A nonzero ordinary integer valuation cannot be divisible by $n^2$
for every $n$.  Thus this explicit toric construction does not itself
supply the full metabelian lift of an infinite-support ordinary pair.
The possibility of other compatible Jacobian coordinates in
(A.68.12)--(A.68.14) remains open.
\end{proof}

